\documentclass[final,3p,times,twocolumn]{elsarticle}
\usepackage{amssymb} % various useful mathematical symbols
\usepackage{newunicodechar}
\newunicodechar{fi}{fi}
\newunicodechar{ff}{ff}
\usepackage{lipsum}
\usepackage{caption}
\usepackage{subcaption}
\usepackage{xcolor}
\usepackage{graphicx}
\usepackage{float}
\usepackage{amsthm} % extended theorem environments
\usepackage{lineno} % line numbers
\graphicspath{{figures/}} % path to folder with figures

\journal{Archives of Acoustics}

\begin{document}

\begin{frontmatter}

\title{Seismic CPT for testing depth stabilization with LC columns}

\author[label1,label2]{Per Lindh}

\affiliation[label1]{organization={Swedish Transport Administration, E-mail: per.lindh@trafikverket.se},
             addressline={Neptunigatan 52, P.O. Box 366},
             city={Malmö},
             postcode={SE-201-23},
	country={Sweden}}
\affiliation[label2]{organization={Lund University, Lunds Tekniska Högskola LTH (Faculty of Engineering), Department of Building and Environmental Technology, Division of Building Materials, E-mail: per.lindh@byggtek.lth.se},
             addressline={Box 118},
             city={Lund},
             postcode={SE-221-00},
             country={Sweden}}

\author[label3]{Polina Lemenkova}
\affiliation[label3]{organization={Université Libre de Bruxelles (ULB), École polytechnique de Bruxelles (EPB, Brussels Faculty of Engineering), Laboratory of Image Synthesis and Analysis (LISA). E-mail: polina.lemenkova@ulb.be},
             addressline={Campus de Solbosch, ULB -- LISA CP165/57, Avenue Franklin D. Roosevelt 50},
             city={Brussels},
             postcode={B-1050},
             country={Belgium}}

\begin{abstract} In this project, seismic (Cone Penetration Testing) CPT has been evaluated as a test method to determine the properties of lime/cement (LC) columns. The reason for using seismic CPT is that there both seismic value obtained in the form of P-wave (compression wave), S-wave (shear wave) and a peak pressure. We found a connection between probing resistance of soil and strength in the stabilized columns. The results of the seismic tests obtained during fieldwork survey were evaluated and compared against seismic values obtained in the laboratory of Swedish Geotechnical Institute (SGI). The CPT was used to evaluate the quality of the LC columns of the reinforced soil. The adjacent projects have been based on the use of bending elements (bender elements) and resonant frequency measurement in the laboratory big project \emph{Seismic testing of stabilized soil with bender elements in triaxial tests} (Lindh, 2016a). This is a simple technique that usually gives an unequivocal answer to the resonance frequency. The homogeneity of the specimen can be examined visually and the relationship between P and S wave shows the quality of the specimen. In the field, this possibility of visually checking the quality does not exist, but here the collected measurement data must be sufficiently good. When it comes to seismic CPT, it has proven to be critical partly to evaluate whether the CPT is inside or outside the pillar. A low tip resistance may mean the probe is off the post or the post is bad. In part, one needs to be able to excite a well-defined P- and/or S-wave for the seismic measurements. Alternative seismic methods such as Up-hole technology e.t.c. should be evaluated.
\end{abstract}

\begin{keyword}

civil engineering \sep stabilization \sep compressive strength \sep soil \sep cement \sep lime \sep seismic waves

\PACS 81.40.Cd \sep 81.40.Ef \sep 62.20.Qp \sep 83.50.Xa \sep 45.70.Mg \sep 92.40.Lg \sep 81.40.Lm \sep 62.20.M-
\MSC 76Axx \sep 74Exx \sep 74Fxx

\end{keyword}

\end{frontmatter}

%% \linenumbers

%% main text
\linenumbers
\section{Introduction}

\subsection{Background}

%% text

\subsection{Objectives and goals}
The background to the project is the need to adequately evaluate the quality of KC columns in terms of strength and homogeneity. The purpose is to better connect laboratory measurements and field measurements to obtain a better and safer optimization of binders and binder types. In this sub-project, the aim is to evaluate seismic 

\section{Methods}

\subsection{Fieldwork}
The area of Norvik consisted of a bay until the beginning of the 1980s, surrounded by a hilly strip of land and a mountainous island. The material parameters of the natural clay are reported in Figure 1 and Figure 2. 
During the construction of a gas storage facility in Nyn\"ashamn, the bay was filled again with explosives from the construction. The filling was carried out by tipping masses from several fronts, which meant that large volumes of clay were enclosed in the blast-stone filling. Two large areas, northern and southern, have been identified, with up to 30 m of clay. A number of geotechnical field investigations in the clay areas have been carried out in the years 1982, 2008, 2010 and 2011. Geomind 2015.

\subsection{Laboratory Tests}

\subsubsection{Choice of examination room}
When choosing an examination room for a test such as the one now carried out, it is desirable that the room is well examined and documented both in terms of field examinations and laboratory examinations. Furthermore, it is desirable that seismic measurements are carried out in the laboratory as a reference to the measurements in the field. In this project, the choice fell on the port of Norvik, where the area was well documented in terms of in situ conditions on natural soil and that the laboratory measurements on stabilized soil had been carried out with measurements of P-wave and flexwave (shear wave). Lindh (2016b).

\subsubsection{Determination of the material parameters of the natural clay}
The natural water ratio, wN, varies between 40 and 100\% and the yield strength, wL, also varies between 40 and 100\%. The density mostly varies between 1.5 to 1.7 tons per cubic meter. The corrected shear strength, cu,korr, is mainly between 5 and 15 kPa and the clay is low to medium sensitive.

\subsubsection{Evaluation of binders}
Testing of binding agent and amount of binding agent was carried out on clay from test excavations carried out at Norvik. To evaluate binder combinations for depth stabilization, an experimental setup called simplex centroid design was used (Montgomery, 1996, Myres and Montgomery, 1995). An experimental set-up according to the simplex centroid with three binders means that the experimental room consists of a triangle with pure binders in the corners and a mixture of two binders along the triangle's stripes. Inside the triangle is a combination of all three binders, the total binder content is constant. Each point on the triangle represents 100\% binder, but the mutual ratio between the binders varies depending on where in the triangle the test point is located. The tests were carried out according to SGI's laboratory quality handbook, document 29a, rev c. Three different binders have been used, cement, lime and lime kiln dust (English; lime kiln dust, LKD). The cement used in the tests has been a CEM II/A and the lime is a burnt lime, named Ql in the figures (from the English Quick lime).

From these three binders, nine different binder combinations have been tested. The mixing quantities corresponded to 80, 100 and 120 kg of binder per cubic meter of clay. The curing times for the samples (cement/lime 50/50) were determined to be 7, 14, 21 and 28 days. Some extra specimens were produced for printing at 80 days. 

The experimental set-up was performed as a three-factor Simplex centroid with a limitation of the LKD content. The restriction meant that the LKD content was maximized to 50\%. The trial surface in a three-factor Simplex centroid is structured as a triangle. The purpose of using this type of statistical trial planning is to minimize the number of trials while maintaining statistical significance and to ensure that both positive and negative interactions between the binder components are detected. 

Each point on the surface corresponds to 100\% binder. The selected trial points are in this case marked with blue circles, see Figure 3. Each trial point was performed as a double trial. The corners of the triangle represent 100\% of only one binder. The center of the triangle corresponds to 33\% of the three different binders. The points on the lower edge of the triangle correspond to either a pure binder, (cf. CEM II or Ql) alternatively a mixture of 50\% CEM II and 50\% Ql. Three different Simplex experiments were performed. These were equivalent to 80, 100 and 120 kg of binder per cubic meter of clay.

For this, extra reference tests were carried out which were to be used for calibration of seismic measurement in relation to compressive strength at 7 and 14 days, see Figure 4.

\subsection{Laboratory seismic measurements}
Natural resonance frequency (fn) (also called natural natural frequency) is the number of oscillations per second (Hz) in a test body that is allowed to oscillate and swing freely without damping. The lowest resonant frequency is called the fundamental mode or fundamental mode. All natural resonance frequencies can be physically related to the elastic constants, E-modulus (E) and transverse contraction number () or seismic velocities such as primary wave (P-wave) and secondary wave (S-wave or shear wave). For specimens with a length twice the diameter ($L/D\geq2$), the one-dimensional (1D) wave propagation velocity can be calculated as Equation \ref{eq:01}

\begin{equation}\label{eq:01}
	V_{P1D} =2Lf_d
\end{equation}

where $f_d$ is the damped resonance frequency (Rydén, 2013).
The seismic measurements were performed as a resonance frequency measurement of P-wave and flex-wave, see Figure 5 and Figure 6.

The advantage of P-wave measurements is that these can be carried out on specimens while they harden in the sleeve without the sleeve affecting the measurements to any great extent. The disadvantage of measuring the P-wave is that it is strongly affected when the soil starts to become saturated with water, when Sr goes towards one. When the earth is saturated with water, not only the earth's P-wave velocity is measured, but also the water's P-wave velocity, which is about 1500 m/s, which means incorrect results. This is usually not a problem in laboratory-made specimens, as these cannot be saturated with water without very high pressures. 

By instead measuring the shear wave when it is not affected by high water saturation levels, this problem is bypassed. The measurement of the true shear wave is, however, more difficult with a free-free resonant column setup. A common, but not entirely accurate, way to evaluate the shear wave is to measure the bending mode also called the flex mode or transverse mode (Verástegui-Flores et al, 2015). According to this article, the shear wave velocity is underestimated by about 5\% for a specimen with a slenderness factor of 2 (=length/diameter) and a = 0.2. The setup for measuring the flex mode is shown in Figure 6.

The evaluation of the shear wave velocity is calculated according to Equation \ref{eq:02}. 

\begin{equation}\label{eq:02}
	V_s = 2Lf_f
\end{equation}

To measure true hear waves in laboratories, you have to use bending elements (so-called bender elements), this methodology is cumbersome and takes a long time. To get around this, a new methodology has been tried at SGI to evaluate the torsional mode and thus be able to better calculate the shear wave, see Figure 7.

However, this methodology is not yet sufficiently tested and must be verified with more measurements but also analyzed using FEM to check that the correct mode is evaluated. To evaluate how both P-wave velocity and shear strength develop with time, reference samples were manufactured. The choice of the P-wave velocity in the laboratory examinations is that it can also be measured on samples stored in sleeves. However, the length of the sleeves must be equal to the length of the sample. Otherwise, the resonance frequency of the sleeve is measured, see Appendix 1. The reference samples consisted of a standard recipe consisting of 50\% cement and 50\% quicklime. The binder quantities were chosen corresponding to 80, 100 and 120 kg/m3. After 7 days, P-wave velocity and then shear strength were evaluated on two samples from each binder quantity. This was repeated at 14, 21, 28 and 80 days. This type of correlation has previously been used for surface-stabilized specimens with good results (Rydén et al, 2006). The choice of 80 days was based on degree days to compare with the 28 day strength of samples stored at 20 oC. A better procedure would have been to use maturity numbers (MT) instead of degree days, see SGI Report 30, j.m.f. section 3.1.2.

\subsection{Seismic CPT}

The CPT probe is performed according to the standard SS-EN ISO 22476-1:2012. During a CPT probing, a cylindrical probe with a cross-sectional area of 1000 mm2 is driven where the probe has a tip angle of 60 degrees. The lowering takes place at a constant speed of 20 mm/s. At the probe tip, the force required to drive the probe down is measured. Through a backlash coupling, the casing friction is measured, which can thus be distinguished from the tip pressure. Furthermore, the pore pressure generated during the drive-down is measured via a filter system just behind the probe tip. Figure 8 shows the construction of a CPT probe.

During the probing, the tip resistance is recorded, which is the force per unit area obtained by dividing the measured force by the cross-sectional area of the tip. This peak pressure is denoted qc if the pore pressure is not taken into account and qt if the value is corrected taking into account error sources from the pore effect. In the special case when $u \approx 0$ or is negligible, qc becomes qt. The tip resistance is given in MPa or kPa. The CPT probe can be reported according to Figure 9.

Seismic CPT involves, in addition to a regular CPT sounding, also a measurement of shear and compression waves. The measurements are performed with a series-connected device that is placed between the CPT probe and the probing rods. The measurements are usually performed by stopping the CPT probe at the level you want to measure at and the measurement is performed. However, there are systems with continuous measurement, but this must be stopped after all when splicing bars.

There are two different principles, one principle is based on two rounds of accelerometers with a fixed distance between the accelerometers, see Figure 10. The accelerometers measure in the x-, y- and z-direction.

This principle is called "true-time" as the calculation of walking speed is based on the time measurement between the first wave propagation to accelerometers 1 and 2. As the distance between the accelerometers is constant, the speed can be calculated directly.
The second principle is based on carrying out the measurements with only one set of accelerometers with registration in three directions (x-, y- and z-direction). The measurement is first performed at one level, after which the CPT probing continues and a new measurement is performed at the next level, see Figure 11.

\subsubsection{Evaluation of CPT}

The evaluation of the CPT probe has been performed using Conrad version 3.1. In a fine-grained soil, the shear strength is primarily evaluated as the net peak pressure. An alternative method for clays is to use the generated pore overpressure. The empirical relationships developed for Swedish conditions are based on a wealth of field and laboratory data from a variety of soil types. In a natural clay soil, the relationship between net peak pressure and shear strength is sensitive to the soil's yield strength wL and to some extent also to the degree of over-consolidation OCR, which is often described according to the relationship in Equation \ref{eq:03}

\begin{equation}\label{eq:03}
	c_u = \frac{q_t -\sigma _{\nu 0}}{13,4+6,65W_L} {\left(\frac{OCR}{1,3}\right)}^{-0,2}
\end{equation}

\begin{figure*}[!htbp]
\centering
	\begin{subfigure}[b]{.5\textwidth}
		\centering
			\includegraphics[width=0.9\textwidth]{AEG35_1}
		\caption{}
	\end{subfigure}%
	\begin{subfigure}[b]{.5\textwidth}
		\centering
			\includegraphics[width=0.9\textwidth]{AEG35_2}
		\caption{}
	\end{subfigure}%
	\\ \vfill \vspace{1mm}
	\begin{subfigure}[b]{.5\textwidth}
		\centering
			\includegraphics[width=0.9\textwidth]{AEG35_3}
		\caption{}
	\end{subfigure}%
	\begin{subfigure}[b]{.5\textwidth}
		\centering
			\includegraphics[width=0.9\textwidth]{AEG35_4}
		\caption{}
	\end{subfigure}
	\begin{subfigure}[b]{.5\textwidth}
		\centering
			\includegraphics[width=0.9\textwidth]{AEG35_5}
		\caption{}
	\end{subfigure}%
	\begin{subfigure}[b]{.5\textwidth}
		\centering
			\includegraphics[width=0.9\textwidth]{AEG35_6}
		\caption{}
	\end{subfigure}%
	\\ \vfill \vspace{1mm}
	\begin{subfigure}[b]{.5\textwidth}
		\centering
			\includegraphics[width=0.9\textwidth]{AEG35_7}
		\caption{}
	\end{subfigure}%
	\begin{subfigure}[b]{.5\textwidth}
		\centering
			\includegraphics[width=0.9\textwidth]{AEG35_8}
		\caption{}
	\end{subfigure}
	\\ \vfill \vspace{1mm}
	\begin{subfigure}[b]{.5\textwidth}
		\centering
			\includegraphics[width=0.9\textwidth]{AEG35_9}
		\caption{}
	\end{subfigure}%
\caption{AEG35}
\label{figAEG35}
\end{figure*}

\section{Results}

\lipsum[1-2]

\newpage
\lipsum[1-10]

\newpage
\section{Discussion}

\lipsum[1-12]

\section{Conclusion}

\lipsum[1-12]

\section{Acknowledgements}
The authors are thankful to the two anonymous reviewers who provided careful reading and commenting of this work. We are grateful  to Helen Åhnberg and Rolf Larsson for project idea development. The authors thank the Swedish Transport Administration, Skanska and SBUF for the financial contribution to this project. The authors gratefully acknowledge Ports of Stockholm, NCC and Geomind for assistance with data collection and support with the surveys. The authors express the gratitude to Cemex and SMA for adhesives and other support related to the project. Roger Wisén (Ramböll) and Nils Rydén (LTH/Peab) are acknowledged support for their support with simulations. Thanks go to Torbjörn Edstam, Skanska for the review and valuable comments.

%% The Appendices part is started with the command \appendix;
%% appendix sections are then done as normal sections
%% \appendix

%\bibliographystyle{elsarticle-num} 
\bibliographystyle{plain}
\bibliography{References_AA.bib}
\nocite*{}

\end{document}
\endinput
